
\documentclass{article}
\usepackage{colortbl}
\usepackage{makecell}
\usepackage{multirow}
\usepackage{supertabular}

\begin{document}

\newcounter{utterance}

\centering \large Interaction Transcript for game `cladder', experiment `full\_v1.5\_default', episode 1133 with reflexion{-}Qwen3.5{-}9B.
\vspace{24pt}

{ \footnotesize  \setcounter{utterance}{1}
\setlength{\tabcolsep}{0pt}
\begin{supertabular}{c@{$\;$}|p{.15\linewidth}@{}p{.15\linewidth}p{.15\linewidth}p{.15\linewidth}p{.15\linewidth}p{.15\linewidth}}
    \# & \multicolumn{2}{c}{Player} && \multicolumn{2}{c}{Game Master} \\
    \hline

    \theutterance \stepcounter{utterance}  
    & & & \multicolumn{4}{p{0.6\linewidth}}{
        \cellcolor[rgb]{0.9,0.9,0.9}{
            \makecell[{{p{\linewidth}}}]{
                \texttt{\tiny{[P1$\langle$GM]}}
                \texttt{You have attempted this same task before and did not succeed. These are the plans you wrote after those attempts.} \\
\texttt{Trial \#1: New plan: I incorrectly concluded "No" because I stopped my reasoning prematurely after calculating the positive Average Treatment Effect (0.3147) and failed to explicitly state the final conclusion that a positive ATE implies treatment increases recovery. Since the question asks "Will receives treatment increase the chance of recovery?", the answer must be "Yes". Next time, I will ensure that after performing the necessary calculations, I explicitly compare the result to zero and provide the definitive "Yes" or "No" answer as the very first word of my response, followed immediately by the supporting evidence.} \\
\\ 
\\ 
\texttt{You are an expert in causal inference. The following question is not a typical commonsense query, but rather a meticulously designed question created by a professor specializing in causal inference, intended to assess the students' mastery of the course content.} \\
\texttt{Imagine a self{-}contained, hypothetical world with only the following conditions, and without any unmentioned factors or causal relationships: Kidney stone size has a direct effect on treatment and recovery. Treatment has a direct effect on recovery. For patients who have small kidney stones and not receiving treatment, the probability of recovery is 3\%. For patients who have small kidney stones and receiving treatment, the probability of recovery is 33\%. For patients who have large kidney stones and not receiving treatment, the probability of recovery is 60\%. For patients who have large kidney stones and receiving treatment, the probability of recovery is 93\%. The overall probability of large kidney stone is 49\%. Will receives treatment increase the chance of recovery?} \\
\texttt{Start your answer with 'Yes' or 'No', followed by additional reasoning or evidence to support your explanation.} \\
            }
        }
    }
     \\ \\

    \theutterance \stepcounter{utterance}  
    & \multicolumn{4}{p{0.6\linewidth}}{
        \cellcolor[rgb]{0.561,0.737,0.561}{
            \makecell[{{p{\linewidth}}}]{
                \texttt{\tiny{[P1$\rangle$GM]}}
                \texttt{Yes, receiving treatment increases the chance of recovery.} \\
\\ 
\texttt{To determine this, we must calculate the Average Treatment Effect (ATE), which represents the expected difference in recovery probability between receiving treatment and not receiving treatment across the entire population. The ATE is defined as \$E{[}Y(1) {-} Y(0){]}\$, where \$Y(1)\$ is the potential outcome with treatment and \$Y(0)\$ is the potential outcome without treatment.} \\
\\ 
\texttt{First, let us identify the potential outcomes for each subgroup based on the provided probabilities:} \\
\texttt{1.  **Small Kidney Stones:**} \\
\texttt{    *   Probability of recovery with treatment (\$P(Y=1|Small, T=1)\$): 33\%} \\
\texttt{    *   Probability of recovery without treatment (\$P(Y=1|Small, T=0)\$): 3\%} \\
\texttt{    *   Treatment Effect for Small Stones: \$33\textbackslash{}\% {-} 3\textbackslash{}\% = 30\textbackslash{}\%\$ (or 0.30)} \\
\\ 
\texttt{2.  **Large Kidney Stones:**} \\
\texttt{    *   Probability of recovery with treatment (\$P(Y=1|Large, T=1)\$): 93\%} \\
\texttt{    *   Probability of recovery without treatment (\$P(Y=1|Large, T=0)\$): 60\%} \\
\texttt{    *   Treatment Effect for Large Stones: \$93\textbackslash{}\% {-} 60\textbackslash{}\% = 33\textbackslash{}\%\$ (or 0.33)} \\
\\ 
\texttt{Next, we calculate the overall ATE by weighting these subgroup effects by the prevalence of each stone size. The problem states that the overall probability of having a large kidney stone is 49\% (\$P(Large) = 0.49\$). Consequently, the probability of having a small kidney stone is \$1 {-} 0.49 = 0.51\$.} \\
\\ 
\texttt{The calculation for the ATE is:} \\
\texttt{\$\$ATE = P(Small) \textbackslash{}times \textbackslash{}text\{Effect\}\_\{Small\} + P(Large) \textbackslash{}times \textbackslash{}text\{Effect\}\_\{Large\}\$\$} \\
\texttt{\$\$ATE = 0.51 \textbackslash{}times 0.30 + 0.49 \textbackslash{}times 0.33\$\$} \\
\texttt{\$\$ATE = 0.153 + 0.1617\$\$} \\
\texttt{\$\$ATE = 0.3147\$\$} \\
\\ 
\texttt{The calculated Average Treatment Effect is 0.3147 (or 31.47\%). Since the ATE is strictly greater than zero, it indicates that, on average, receiving treatment increases the probability of recovery by approximately 31.5 percentage points compared to not receiving treatment. Therefore, the answer is definitively yes.} \\
            }
        }
    }
    & & \\ \\

    \theutterance \stepcounter{utterance}  
    & & & \multicolumn{2}{p{0.3\linewidth}}{
        \cellcolor[rgb]{0.9,0.9,0.9}{
            \makecell[{{p{\linewidth}}}]{
                \texttt{\tiny{[GM$|$GM]}}
                \texttt{Yes} \\
            }
        }
    }
    & & \\ \\

    \theutterance \stepcounter{utterance}  
    & & & \multicolumn{2}{p{0.3\linewidth}}{
        \cellcolor[rgb]{0.9,0.9,0.9}{
            \makecell[{{p{\linewidth}}}]{
                \texttt{\tiny{[GM$|$GM]}}
                \texttt{yes} \\
            }
        }
    }
    & & \\ \\

    \theutterance \stepcounter{utterance}  
    & & & \multicolumn{2}{p{0.3\linewidth}}{
        \cellcolor[rgb]{0.9,0.9,0.9}{
            \makecell[{{p{\linewidth}}}]{
                \texttt{\tiny{[GM$|$GM]}}
                \texttt{game\_result = WIN} \\
            }
        }
    }
    & & \\ \\

\end{supertabular}
}

\end{document}
